\documentclass[UTF8,12pt,a4paper]{ctexart}
\usepackage{amsmath,amssymb,amsthm}
\usepackage{booktabs}
\usepackage{graphicx}
\usepackage{algorithm}
\usepackage{algpseudocode}
\usepackage{geometry}
\usepackage{hyperref}
\geometry{margin=2.5cm}

\newtheorem{definition}{定义}
\newtheorem{assumption}{假设}
\newtheorem{proposition}{命题}
\newcommand{\Cmax}{C_{\max}}

\title{闭环优于一次成型：FJSP+MAPF耦合系统的\\在线重调度编排与承诺一致性修订}
\author{SkyResearch}
\date{2026年9月1日 \quad v1 草稿}

\begin{document}
\maketitle

\begin{abstract}
在 FJSP 与 MAPF 耦合的车间系统中（I-FJSP-MAPF，见姊妹篇规约），现有最先进做法
（如 Li 等 2026 的 JSP+MAPF 框架）采用``一次规划—执行期修正''的准开环范式。
本文（1）将耦合问题形式化为\textbf{在线}版本 O-IFJSP-MAPF：外生扰动过程、
状态轨迹、闭环策略空间与扰动遗憾（disruption regret）；（2）定义\textbf{承诺
一致性修订}语义——任何在线重规划必须与已执行前缀一致且不违背已冻结承诺，
该语义与引擎的 PlanRevision/激活机制一一对应；（3）提出\textbf{恢复编排器}
（Recovery Orchestrator）：以 $(\tau,\kappa)$ 策略空间补上引擎缺失的闭环
策略层；（4）40 个扰动 episode 的试点给出三个出人意料的结论：
一次成型 CP-SAT 计划对 200 步单机停机仅退化 4--14\%（``等待—吸收''的
被动鲁棒性）；事件触发的全域修订在在制扰动下\emph{全部}激活失败
（未解决承诺）而局部修复多为空作用域——\textbf{闭环的瓶颈在承诺解除
机制而非触发策略}；失败的重规划尝试本身有 7\% 的负作用。该负结果把
闭环协同的研究议程从策略层重定向到机制层（可违约承诺语义）。
\end{abstract}

\section{引言}

姊妹篇（方向2）定义了 I-FJSP-MAPF 的静态规约与分层执行语义
$\mathcal{H}=(\sigma_J,\sigma_A,\sigma_R)$。本篇处理其在线版本：
执行期内发生外生扰动 $\xi$（机器故障与修复、AGV 故障、临时路障）时，
\emph{谁}在\emph{何时}以\emph{多大范围}重调度？

现状的三个断层：
\begin{enumerate}
  \item \textbf{引擎断层}：SkyEngine（及同类系统）的调度层已暴露
    修订 API（残差问题构建、承诺冻结、激活与回滚），但\emph{不存在}
    调用它们的策略组件——扰动下系统实际上是开环的；
  \item \textbf{文献断层}：动态 FJSP 工作考虑机器/AGV 故障但不含栅格
    冲突路由；MAPF 的滚动方法（RHCR~\cite{li2021rhcr}）考虑执行偏差
    但无调度层；Li 等~\cite{li2026framework} 自述其框架为一次性
    ``generate-then-refine''，闭环列为未来工作；
  \item \textbf{评测断层}：无同时度量扰动遗憾与重调度代价
    （修订次数、求解时间、计划紧张度）的协议。
\end{enumerate}

本文贡献：O-IFJSP-MAPF 在线规约（\S\ref{sec:form}）、承诺一致性修订语义
（\S\ref{sec:revision}）、恢复编排器设计与策略空间（\S\ref{sec:orch}）、
四策略扰动试点（\S\ref{sec:pilot}）。

\section{在线问题规约 O-IFJSP-MAPF}\label{sec:form}

\subsection{外生扰动过程}

在静态输入（姊妹篇 \S3.1）之上，定义扰动过程
$\xi=\{\xi_t\}_{t\ge0}$，$\xi_t$ 为 $t$ 时刻发生的事件集合，取值于：
\begin{align}
\xi_t \in \{&\texttt{m\_down}(m,d),\ \texttt{m\_up}(m),\
\texttt{a\_down}(k,d),\ \texttt{a\_up}(k),\ \texttt{obs}(v,d)\}
\end{align}
分别表示机器 $m$ 故障 $d$ 步/修复、AGV $k$ 故障/修复、单元 $v$ 临时封闭
$d$ 步。扰动可为预定表（schedule）或随机过程（MTBF/泊松到达），
由外生种子独立于决策种子控制（可复现性）。

\subsection{状态、策略与目标}

\begin{definition}[在线耦合调度策略]
策略 $\pi$ 在每步 $t$ 观测状态
$s_t=(\text{剩余工序},\ \text{机器/AGV 状态},\ \text{AGV 位置},\ \text{在途任务},\ \text{当前计划}\Pi_t,\ \xi_{\le t})$，
输出（i）是否触发修订及修订范围（本篇的编排决策），（ii）运输任务分派，
（iii）AGV 动作。轨迹由分层求解器 $\mathcal{H}$ 与环境演化共同产生。
\end{definition}

\begin{definition}[扰动遗憾]
固定扰动场景 $S$（含空场景 $S_0=\varnothing$），定义
\[
\mathrm{Regret}(\pi;S) \;=\; \Cmax(\pi,S)\;-\;\Cmax(\pi,S_0),
\qquad
\mathrm{RRegret}(\pi;S)=\frac{\Cmax(\pi,S)}{\Cmax(\pi,S_0)}.
\]
\end{definition}

目标：在场景分布 $\mathcal{D}$ 上最小化
$\mathbb{E}_{S\sim\mathcal{D}}[\mathrm{RRegret}]$，同时以修订次数
$N_{\text{rev}}$ 与编排求解时间 $T_{\text{plan}}$ 为约束/次级目标。
（竞争比式分析留作后续理论工作。）

\section{承诺一致性修订}\label{sec:revision}

在线重规划不能从零开始：执行已发生、承诺已发出。

\begin{definition}[可行修订]
$\Pi_{t+1}=\rho(\Pi_t\mid s_t)$ 为可行修订，若满足：
\begin{enumerate}
  \item \textbf{执行一致}（weak execution consistency）：$\Pi_{t+1}$ 在
    $t$ 时刻前的调度/路径与已执行轨迹逐点一致；
  \item \textbf{承诺一致}（commitment consistency）：已开工工序不被迁移、
    在途运输任务（AGV 已取货）必须按原目的地交付、已冻结指派不变；
  \item \textbf{残差最优}：在 1--2 约束下，$\Pi_{t+1}$ 为残差问题
    （剩余工序 + 未冻结决策）上的（近似）最优。
\end{enumerate}
\end{definition}

\begin{proposition}[锚定下界不变性]
若修订保持承诺一致性，则任意时刻的已完成+在途部分构成最终 $\Cmax$ 的
单调不减下界；因此 $\mathrm{Regret}\ge 0$ 且仅由未冻结部分的次优性与
等待代价构成。
\end{proposition}
（证明平凡，从承诺一致性直接得出；意义在于把遗憾归因拆分为
``冻结代价 + 残差再优化不足''两部分，分别对应 $\kappa=\texttt{partial}$
冻结过多与 $\kappa=\texttt{full}$ 求解限时不足两种失效模式。）

\paragraph{与引擎的对应。}残差问题构建 = \texttt{build\_residual\_problem}；
承诺分类（transport/processing 冻结）= \texttt{frozen\_operations}；
激活校验 = \texttt{execution\_ready} 与激活失败回滚；
修订账本与紧张度 = \texttt{PlanRevisionLedger} 与
\texttt{ScheduleRevisionTracker}。本篇不改动这些机制，只补策略层。

\section{恢复编排器}\label{sec:orch}

\begin{algorithm}[h]
\caption{Recovery Orchestrator（每步调用）}
\begin{algorithmic}[1]
\State $E_t \gets$ \texttt{injector.get\_step\_events()}
\If{$\tau=\texttt{never}$ 或 $\sigma_J$ 为规则求解器} \State \textbf{return} \EndIf
\State $e \gets \textsc{Trigger}(\tau, E_t, t;\ \Delta_{\min})$ \Comment{含最小重规划间隔 $\Delta_{\min}$ 防风暴}
\If{$e=\bot$} \State \textbf{return} \EndIf
\If{$\kappa=\texttt{partial}$ \textbf{and} $e$ 为机器故障}
  \State $\rho \gets$ \texttt{request\_partial\_schedule\_repair}（围绕故障机器的影响域）
\ElsIf{$\kappa=\texttt{full}$}
  \State $\rho \gets$ \texttt{request\_plan\_revision}（全域残差, 承诺冻结, 激活）
\EndIf
\State 激活失败 $\Rightarrow$ 记账回滚（开环降级一步）
\end{algorithmic}
\end{algorithm}

策略空间与对照：
\begin{itemize}
  \item \textbf{greedy-reactive}：规则调度每步重建（规则级闭环，无修订账本）；
  \item \textbf{cpsat-static}：$\tau=\texttt{never}$（一次成型开环对照）；
  \item \textbf{cpsat-full}：$\tau=\texttt{event},\ \kappa=\texttt{full}$；
  \item \textbf{cpsat-partial}：$\tau=\texttt{event},\ \kappa=\texttt{partial}$。
\end{itemize}
假设：H1 $\texttt{cpsat-full/partial}$ 的 RRegret 显著低于
$\texttt{cpsat-static}$（闭环价值）；H2 扰动强度增大时
$\texttt{partial}$ 的求解开销优势扩大（范围效应）；H3
$\texttt{static}$ 在无扰动场景最优（开环无修订税）。

\section{试点实验}\label{sec:pilot}

\textbf{设置}：$\{$mk01, mk02$\}$ $\times$ 迷宫地图 $\times$ 4 AGV $\times$
场景 $\{S_0, S_1{=}\texttt{m\_down}@150{\times}40,\ 
S_2{=}S_1{+}\texttt{a\_down}@300{\times}20,\ 
S_3{=}\texttt{mild\_stochastic},\ S_4{=}\texttt{m\_down}@80{\times}200\ (\text{强})\}$
$\times$ 4 策略；greedy 系用内置 A$^*$ 路由，CP-SAT 系用滚动 EECBS；
episode 上限 4096 步，进程级硬超时。

\begin{table}[h]
\centering
\begin{tabular}{llrrrrr}
\toprule
实例 & 策略 & $S_0$ & $S_1$ & $S_2$ & $S_3$ & $S_4$(强) \\
\midrule
mk01 & greedy-reactive & 1099 & 1099 & 1155 & 977 & \textbf{595} \\
mk01 & cpsat-static    & \textbf{408} & \textbf{408} & \textbf{423} & \textbf{449} & \textbf{467} \\
mk01 & cpsat-full      & 408 & 408 & 423 & 449 & 467 \\
mk01 & cpsat-partial   & 408 & 408 & 423 & 449 & 467 \\
\midrule
mk02 & greedy-reactive & 4096$^\dagger$ & 3417 & 781 & 1103 & 1132 \\
mk02 & cpsat-static    & \textbf{410} & \textbf{410} & \textbf{426} & \textbf{427} & \textbf{426} \\
mk02 & cpsat-full      & 410 & 410 & 426 & 427 & 457$^\ddagger$ \\
mk02 & cpsat-partial   & 410 & 410 & 426 & 427 & 426 \\
\bottomrule
\end{tabular}
\caption{闭环试点 makespan（步）。$\dagger$：饥饿型活性失稳（姊妹篇发现6，
4096 步未完工）；$\ddagger$：失败修订尝试的规划延迟副作用。
三个 CP-SAT 策略几乎处处同分——见发现 F2。}
\label{tab:closemain}
\end{table}

\textbf{发现}：
\begin{enumerate}
  \item \textbf{F1（H3 成立）}：无扰动时名义开环 CP-SAT 大幅最优
    （408 vs 1099；mk02 上 greedy 甚至失稳）；即使 $S_4$ 的 200 步
    单机停机也只使静态计划退化 4--14\%——运输请求的``按计划 ready
    时刻释放 + 机器等待''语义构成了\emph{被动鲁棒性}；
  \item \textbf{F2（H1 在当前栈被拒绝——机制瓶颈）}：事件触发的
    全域修订（U4）在在制扰动下\emph{全部}激活失败
    （``未解决承诺''：被打断的在制工序使残差问题不可满足承诺约束），
    局部修复（U3）多数返回空作用域（计划松弛吸收扰动）——因此
    full/partial/static 三臂不可区分。\textbf{闭环的瓶颈不在触发策略
    $(\tau,\kappa)$，而在承诺解除机制}：需要``可违约+罚金''的承诺语义
    （引擎侧改造项），这把议程从策略层重定向到机制层；
  \item \textbf{F3}：greedy-reactive 对扰动的\emph{增量}很小
    （mk01: 1099$\to$1155；$S_4$ 下甚至 595 $<$ 1099——停机期间
    负载下降），但其\emph{绝对水平}始终劣于 CP-SAT 系 44--72\%，
    且存在失稳风险——``反应式=鲁棒''的直觉只对增量成立；
  \item \textbf{F4}：失败的重规划尝试有真实代价（mk02 $S_4$：
    full 457 vs static 426，规划延迟挤占执行）——重规划要么激活
    成功，否则不如不做，强化了 F2 的机制先行结论。
\end{enumerate}

\section{正式实验计划（待评估后执行）}

（i）引擎侧：可违约承诺语义（罚金参数化）+ 机器恢复事件的
修订窗口；（ii）实验侧：实例扩至 mk01--mk10 $\times$ 2 地图 $\times$
$K\in\{2,4,6\}$ $\times$ 场景 $\{S_0..S_4,$ stress$\}$ $\times$ 4 策略
$\times$ 3 种子 $\approx$ 1680 episode，约 7--10 小时分批；
（iii）$\tau=$periodic-$K$ 扫描与 $\Delta_{\min}$ 消融。
[FULL\_RUN\_PLACEHOLDER]

\section{结论}

本篇把 FJSP+MAPF 耦合系统从``一次规划''推进到``闭环编排''：给出在线规约、
与引擎修订机制一一对应的承诺一致性语义，以及 $(\tau,\kappa)$ 恢复编排器。
试点量化闭环相对开环的扰动遗憾收益与修订税的权衡。代码见
\texttt{sky\_research} 仓库 \texttt{0901closeloop} 分支。

\bibliographystyle{plain}
\bibliography{references_closeloop}

\end{document}
